% Source: https://artofproblemsolving.com/community/c3h3043107_making_an_arithmetic_sequence
% Stable identity: aops-topic-3043107-q01
The six sums are $a,b,c,a+b,a+c,b+c$. The smallest two are $a,b$, so $b=a+d$. The largest two are $a+c,b+c$. The middle two are $c,a+b$ in one of two orders. If $c<a+b$, the progression forces $c=a+2d$ and $a+b=a+3d$. Substituting $b=a+d$ gives $a=2d$, hence $A=\{0,2d,3d,4d\}$. If $a+b<c$, the third term $a+b$ equals $a+2d$, giving $a=d$; the fourth then gives $c=4d$, so $A=\{0,d,2d,4d\}$. Both sets work, yielding respectively the sums $2d,\ldots,7d$ and $d,\ldots,6d$. These are all possibilities.
